\section{Backward information as a proposal, not a replacement score}
\label{sec:backward-guidance}

The proposal developed here keeps one LEDH particle cloud, the existing
per-particle UKF covariance recursion, OT, Contract--E--Chol and the GenUT
dual-cap correction. Future observations inform where to propose particles;
they are not appended to the observation likelihood. The returned score is
the analytical total derivative of the returned finite log normalizer. This
choice differs from the two-cloud Fisher-score construction: a Fisher
estimate of the model score need not differentiate a finite particle
likelihood. Neither construction has established lower model-score error.

\subsection{What travels backward}

Suppose the whole observation sequence is already available. Write
$f_\theta(x\mid z)$ for the transition density and $g_t(x)$ for the
observation density at time $t$. All densities in the following proposition
are with respect to fixed dominating measures, and the conditioning
normalizers are positive and finite. Future information is the function
$h_t(x)=p_\theta(y_{t+1:T}\mid X_t=x)$, not a sampled trajectory.

\begin{proposition}[Backward information and the smoothing proposal]
\label{prop:bg-backward}
With $h_T=1$,
\begin{align}
 h_{t-1}(z)&=\int f_\theta(x\mid z)g_t(x)h_t(x)\,\dd x,
 \label{eq:bg-backward}\\
 p_\theta(x_t\mid x_{t-1}=z,y_{t:T})
 &=\frac{f_\theta(x_t\mid z)g_t(x_t)h_t(x_t)}{h_{t-1}(z)}.
 \label{eq:bg-smoothing-proposal}
\end{align}
If $h_t$ is constant, it has no effect on the normalized proposal.
\end{proposition}
\begin{proof}
Condition on $X_t=x$ in $p(y_{t:T}\mid X_{t-1}=z)$ and use the Markov
property and conditional independence of observations. This gives the first
identity. Bayes' rule gives the second; a constant factor cancels between
numerator and denominator.
\end{proof}
This is the state-space backward-information argument of
\citet[eqs.~(9)--(10), Proposition~4.1]{heng2020controlled}. Their approximate
dynamic programming fits backward policy corrections using pilot
parent--child pairs. It is not the same operation as following surviving
ancestral trajectories. A fitted function may supply useful future
information without reconstructing any full smoothing path.

\subsection{A precisely defined, frozen-guide proposal}

We first restrict the proposed implementation to full-rank additive Gaussian
transitions $f_\theta(x\mid z)=\mathcal N(x;a,Q)$, where
$a=m_\theta(z)$ and $Q=Q_\theta\succ0$. This excludes singular transitions.
A pilot produces a fixed schedule of guide coefficients $\eta$ before the
evaluated finite program starts. These coefficients, any fit reference
parameter, the observation sequence and base Gaussian draws are fixed when
that program is differentiated. Re-fitting at each evaluated parameter and
then ignoring the fit derivative would compute a different quantity.

For each incoming particle, let a guide evaluator return symmetric
$H\succeq0$ and $b$, possibly depending smoothly on its current anchor $a$.
At this anchor the local proposal factor is
$\exp(-x^\top Hx/2+b^\top x+c)$. The additive constant $c$ need not be
known because it cancels on normalization. There is no extra observation
factor in the importance numerator.

\begin{proposition}[Gaussian proposal and its total tangent]
\label{prop:bg-gaussian}
The normalized product of the Gaussian transition and this local factor is
$q^0=\mathcal N(m,V)$, where
\begin{align}
 P&=Q^{-1},& K&=P+H,& V&=K^{-1},& m&=V(Pa+b).
 \label{eq:bg-gaussian}
\end{align}
For any parameter direction, dots denote total directional derivatives,
including dependence through the incoming state. Then
\begin{align}
 \dot P&=-P\dot QP,& \dot K&=\dot P+\dot H,&
 \dot V&=-V\dot K V,\nonumber\\
 \dot m&=V(\dot P a+P\dot a+\dot b-\dot K m).
 \label{eq:bg-gaussian-tangent}
\end{align}
For fixed $\xi\sim\mathcal N(0,I)$, put $U=m+L\xi$ with
$L=\operatorname{chol}(V)$. Its tangent is
$\dot U=\dot m+\dot L\xi$, where
$\dot L=L\Phi(L^{-1}\dot V L^{-\top})$ and $\Phi$ retains the lower
triangle, halves the diagonal and sets the upper triangle to zero.
\end{proposition}
\begin{proof}
Expand the Gaussian exponent, collect its quadratic coefficient $P+H$ and
linear coefficient $Pa+b$, and complete the square. Since $P\succ0$ and
$H\succeq0$, $K\succ0$; thus $Q,K,V,L$ are invertible. Explicitly,
\begin{align*}
 \dot P Q+P\dot Q&=0,\\
 \dot K V+K\dot V&=0,\\
 \dot K m+K\dot m&=\dot P a+P\dot a+\dot b.
\end{align*}
Multiplication by the appropriate inverse gives each displayed tangent.
Differentiating $LL^\top=V$
and imposing lower triangular $\dot L$ gives the stated Cholesky formula.
The base draw is constant, giving the last identity.
\end{proof}
For a controlled-SMC quadratic policy
$\psi(x)=\exp(-x^\top A x-\beta^\top x-c)$, the adapter is
$H=2A$, $b=-\beta$ when $A\succeq0$. The sign and factor of two agree
with the released Lorenz implementation's shifted Gaussian precision. This
adapter borrows a proposal factor; it does not claim to execute the complete
controlled-SMC twisted-potential algorithm. Policies depending on the current
observation may still be corrected as proposals, but only a future-only
policy can support the specific lookahead interpretation above.

\subsection{What a kernel mixture can contribute}

A separate optional representation is a positive frozen Gaussian mixture
$\widehat h(x)=\sum_k\alpha_k\mathcal N(x;c_k,S_k)$ with
$\alpha_k>0$, $\sum_k\alpha_k=1$, and $S_k\succ0$. Multiplication of
$\widehat h$ by a positive constant has no effect on the following guide.
A normalized mixture is only a hypothesis for the shape of backward
information; backward information need not itself be integrable as a
function of state. No claim of universal representability is made.
Unlike sampling a component label, evaluating this mixture introduces no
categorical discontinuity. Its use here is a new local proposal construction,
not the IWSG estimator of \citet{younis2023}.

\begin{proposition}[Mixture-informed local Gaussian guidance]
\label{prop:bg-mixture}
At the anchor $a$, define $\Lambda_k=S_k^{-1}$ and
\begin{align}
 r_k&=\frac{\alpha_k\mathcal N(a;c_k,S_k)}{\widehat h(a)},&
 v_k&=\Lambda_k(c_k-a),& \bar v&=\sum_k r_kv_k,\nonumber\\
 H&=\sum_k r_k\Lambda_k,& b&=\sum_k r_k\Lambda_kc_k.
 \label{eq:bg-mixture}
\end{align}
Then $H\succ0$ and the local quadratic has the same spatial gradient as
$\log\widehat h$ at $a$, since $b-Ha=\bar v$.
For fixed mixture coefficients,
\begin{align}
 \dot r_k&=r_k(v_k-\bar v)^\top\dot a,&
 \dot H&=\sum_k\dot r_k\Lambda_k,&
 \dot b&=\sum_k\dot r_k\Lambda_kc_k.
 \label{eq:bg-mixture-tangent}
\end{align}
\end{proposition}
\begin{proof}
The gradient of a Gaussian log density is $v_k$ at $a$; differentiation of
the log-sum-exp gives $\bar v$. In particular,
$\dot{\log r_k}=v_k^\top\dot a-\bar v^\top\dot a$ and
$\dot r_k=r_k\dot{\log r_k}$, which gives the claimed tangent.
Differentiate the finite sums to obtain $\dot H$ and $\dot b$. A positive
weighted sum of positive-definite matrices is positive definite, and
substitution gives $b-Ha=\sum_k r_k\Lambda_k(c_k-a)$.
\end{proof}
This is gradient matching, not Hessian matching:
$\nabla^2\log\widehat h(a)=-H+\operatorname{Cov}_r(v_k)$.
Discarding the covariance term defines the proposal approximation explicitly;
it does not justify treating a multimodal mixture as log concave. Nor does a
local Gaussian retain all mixture modes. The omission is not a hidden
stabilizer for the physical likelihood; the proposal density is corrected
and the change in finite-particle behavior still requires evaluation.

\subsection{The corrected finite likelihood and its derivative}

Generate the pre-flow particle from $q^0$, not from the unmodified
transition. Apply the existing LEDH flow using its physical prediction and
current observation. Write its output as $X=F(U)$ and the log-Jacobian
quantity computed by the existing flow as $j$. Its incoming normalized
weight is $W$. The new pre-reset log weight is defined by
\begin{equation}
 \lambda=\log W+\log f_\theta(X\mid z)+\log g_t(X)
                 +j-\log q^0(U).
 \label{eq:bg-weight}
\end{equation}
In particular, subtracting $\log f_\theta(U\mid z)$ instead of
$\log q^0(U)$ is wrong for this proposal. The guide is not multiplied into
$g_t$. Initialization is unchanged: this bounded implementation does not
twist the initial law.

\begin{proposition}[Conditional importance identity and its limitation]
\label{prop:bg-correction}
Condition on the incoming cloud and guide. If $F$ is a differentiable
bijection on the proposal support and $j=\log|\det D_UF|$ is its actual
Jacobian, then for integrable $\varphi$,
\[
 \mathbb E_{U\sim q^0}\!\left[
 \varphi(F(U))\frac{f_\theta(F(U)\mid z)g_t(F(U))e^j}{q^0(U)}\right]
 =\int\varphi(x)f_\theta(x\mid z)g_t(x)\,\dd x.
\]
\end{proposition}
\begin{proof}
Cancel $q^0$ against the sampling density and substitute $x=F(u)$.
The absolute Jacobian converts $\dd u$ into $\dd x$.
\end{proof}
The bijection and actual-Jacobian conditions are substantive. If a localized
flow records only a frozen-coefficient affine determinant while its
coefficients depend on the transported particle, this proposition cannot
certify that implementation. Differentiating the recorded determinant
correctly does not repair that measure mismatch. This work preserves the
existing finite LEDH calculation and audits its derivative; it does not
silently certify those independent flow assumptions. Furthermore, the
deterministic OT/Contract-E/GenUT reset is not multinomial resampling and
does not inherit an unbiased full-filter normalizer from this one-step
identity.

\begin{proposition}[Total derivative of the retained finite program]
\label{prop:bg-total}
Fix base draws, guide coefficients, iteration counts, residual design and
all non-parameter controls. On an open neighborhood where each executed
operation is differentiable, Cholesky factors exist and branch choices are
unchanged, define
\[
 \widehat\ell_\eta(\theta;\xi)
  =\sum_{t=1}^T \log\sum_i\exp(\lambda_t^i),\qquad
 w_t^i=\frac{\exp(\lambda_t^i)}{\sum_j\exp(\lambda_t^j)}.
\]
Propagating all input tangents gives
\begin{align}
 \dot\lambda&=\dot{\log W}+\dot{\log f}+\dot{\log g}
                         +\dot j-\dot{\log q^0},\nonumber\\
 \dot{\widehat\ell}_\eta&=\sum_t\sum_iw_t^i\dot\lambda_t^i,&
 \dot w_t^i&=w_t^i\left(\dot\lambda_t^i-\sum_jw_t^j\dot\lambda_t^j\right).
 \label{eq:bg-total}
\end{align}
This is the total directional derivative of this same finite scalar.
\end{proposition}
\begin{proof}
For $q^0=\mathcal N(m,V)$, put $e=U-m$. Direct differentiation gives
\[
 \dot{\log q^0(U)}=-\tfrac12\operatorname{tr}(V^{-1}\dot V)
 -(V^{-1}e)^\top(\dot U-\dot m)
 +\tfrac12e^\top V^{-1}\dot V V^{-1}e.
\]
Use Proposition~\ref{prop:bg-gaussian} for all its arguments and, when
applicable, Proposition~\ref{prop:bg-mixture} for the moving anchor. The
chain rule gives $\dot\lambda$ and differentiation of log-sum-exp gives
the remaining expressions. Induct over time, carrying state and covariance
tangents through the existing UKF, flow and reset. In the reset, if
$Y=\mathcal T(X,w)$ and $E=\mathcal R(X,w,Y)$, its tangent is
$\mathcal R_X\dot X+\mathcal R_w\dot w+
\mathcal R_Y(\mathcal T_X\dot X+\mathcal T_w\dot w)$, with any declared
ridge/design derivatives as well. Differentiating only $Y$ omits direct
source-moment terms. Apply the same chain rule to GenUT and both caps.
This completes the induction, assuming the invoked component tangents are
correct on that neighborhood.
\end{proof}
The proposition does not equate this finite derivative with the exact model
score. It also does not differentiate the pilot fit: if $\eta=\eta(\theta)$
is recomputed, the omitted term is
$(\partial\widehat\ell/\partial\eta)\dot\eta$.

\begin{proposition}[No automatic score improvement]
\label{prop:bg-no-ranking}
Even exact sampling from a posterior, with constant importance weights,
does not imply minimal variance of a posterior score-functional estimate.
\end{proposition}
\begin{proof}
On three states let the posterior be uniform and let its complete-data
score functional be $(-1,0,1)$. Its expectation is zero. The asymptotic
variance of a self-normalized importance mean is
$\sum_i\pi_i^2 H_i^2/q_i$. Choosing $q=\pi$ gives $2/3$; choosing
$q=(0.45,0.10,0.45)$ gives $40/81<2/3$, despite nonconstant weights.
These scores can be realized by the exponential family proportional to
$\exp(\theta H_i)$ at $\theta=0$. Thus even an ideal distributional
guide does not establish optimal score error. No ranking of finite LEDH
derivatives follows from likelihood ESS.
\end{proof}

\subsection{Executable order and scientific scope}

The bounded implementation is an optional single-cloud proposal extension,
not an alternative LEDH filter. Its order is:
\begin{enumerate}[nosep]
\item Freeze a validated time-indexed guide. Its last future-only factor
      is constant. Fix the existing randomness and numerical controls.
\item Predict physical UKF moments and compute each transition anchor and
      its total tangent from the current reset state.
\item Evaluate $H,b,\dot H,\dot b$; generate $U,\dot U$ from the guided
      Gaussian. With guidance disabled use the original transition draw.
\item Run the shared LEDH flow and its analytical tangent. Evaluate the
      unchanged physical transition and observation densities and the
      actual guided Gaussian denominator in \eqref{eq:bg-weight}.
\item Accumulate the finite log normalizer and derivative; update physical
      UKF covariances. Execute the existing OT, Contract E and GenUT
      dual-cap reset with all tangents. These particles are the next parents.
\end{enumerate}
An OT barycentre has no unique sampled ancestor. Neither its transport
matrix nor the coefficients of a kernel mixture are thereby a model
backward-smoothing kernel. No ancestral tracing or Fisher-score substitution
is part of this implementation.

The first evidence is engineering evidence: disabled-guide parity,
Gaussian square-completion checks, moving-anchor mixture tangents, and
finite differences of the complete same-value recursion. The development
comparator uses $\epsilon=1$ and five Sinkhorn iterations; balancing counts,
ridge, correction and cap settings are reported independently. These are
untuned settings. CPU-only small reference tests do not establish GPU/XLA
readiness. Pilot fitting on actual models, equal-compute repeated-seed score
comparisons, per-scope tuning, batch-native consumer coverage and any
scientific or default promotion remain subsequent work. In particular,
accepting coefficient arrays is not evidence that the replicated fitting
endpoints have been integrated or that their output represents $h_t$.
